\documentclass[PhD, UKenglish]{scrbook}

% Load the checked-out thesis package directly so editor LaTeX builds also
% find it without the TEXINPUTS setting supplied by docs/makefile.
\usepackage[
  twoside=false,
  biblatex=false,
  newtx=false,
  siunitx=false,
  numerr=false,
  numpmcorr=false,
  hepparticles=false,
  mhchem=false
]{../libs/ubonn-thesis/ubonn-thesis}

\begin{document}

\raggedbottom

\frontmatter
\pagestyle{scrplain}

\tableofcontents

\mainmatter
\pagestyle{scrheadings}

\chapter{Introduction}
\section{Motivation and contributions}

\begin{itemize}
  \item Background explaining how to derive a data converter specific to a radiation imaging system (pen-and-paper math).
  \item A set of behavioral models and SPICE simulations to map out the optimization space for SAR ADCs which meet the needs of compact, high-frame-rate image sensors.
  \item A framework/pipeline for translating that into a fabricatable GDS layout, using mostly open-source tools, across a range of PDKs.
  \item Two silicon prototypes built with the simulations and pipeline above.
  \item A hardware DAQ system (COB PCB carrier plus FPGA system, plus control code), with scans to measure all ADC parameters.
  \item Finally, an analysis of measurements from the first generation.
\end{itemize}

\chapter{Background}

\section{Radiation imaging systems}
\subsection{Imaging modalities and detector requirements}
\subsubsection{Stroboscopic high-fluence imaging and high-frame-rate integration}
\subsubsection{Scanning transmission electron microscopy (STEM) and high-frame-rate integration}
\subsection{Electron and X-ray signals and sensor readout}
\subsection{Monolithic and hybrid sensor arrangements}
\subsubsection{Sensor-to-readout connections and pixel pitch}
\subsubsection{Available converter area and power density}
\subsection{Pixel-parallel and column-parallel ADC arrangements}
\subsubsection{Per-pixel conversion versus shared column conversion}
\subsubsection{Multiplexing, aggregate throughput, and frame rate}
\subsection{Single-bit counting and integrating readout}
\subsubsection{Counting dead time, pileup, and false hits}
\subsubsection{Fluence limits of single-bit counting versus integrating readout}
\subsubsection{Amplitude information from multi-bit readout}
\subsection{Imaging rate, dynamic range, and readout requirements}
\subsubsection{Frame rate and aggregate conversion throughput}
\subsubsection{Noise, power density, and pixel-area budgets}

\clearpage

\section{Multi-bit quantization (ADCs)}
\subsection{ADC operation and resolution}
\subsubsection{Quantization, noise, and effective number of bits}
\subsubsection{Static linearity, dynamic range, and distortion}
\subsection{Area, power, and speed performance}
\subsubsection{Conversion rate, latency, and channel density}
\subsubsection{Energy per conversion and resolution-dependent figures of merit}
\subsection{ADC architectures for high-speed imaging}
\subsubsection{Single-slope, flash, cyclic, pipeline, and SAR converters}
\subsubsection{Resolution, area, power, and speed tradeoffs}
\subsection{Fitness of SAR architecture for high-rate imaging}
\subsubsection{Fit to the imaging throughput and per-channel budgets}
\subsubsection{Sampling, CDAC, comparator and sequential-decision limits}
\clearpage

\section{Review of existing designs}
\subsection{High-speed imaging converters in the literature}
\subsection{Development history within the group}
\subsection{Sensor interfaces and converter architectures across designs}
\subsection{Resolution, speed, area, and power comparison}
\subsection{Literature review and position of this work}
\subsubsection{Resolution and conversion-rate density}
\subsubsection{Energy efficiency and area at the target sampling rate}

\chapter{Design of a compact SAR ADC}

\section{Study of a sampling switch}
\subsection{On-resistance and acquisition bandwidth}
\subsubsection{Device size and input-voltage dependence}
\subsubsection{Simple switch versus bootstrapped switch}
\subsection{Sampling noise}
\subsubsection{Thermal-noise and capacitance limits}
\subsubsection{Extracted held-level noise and comparator contributions}
\subsection{Dynamic sampling error}
\subsubsection{Settling error versus sampling-window duration}
\subsubsection{Input bandwidth and effective-resolution limits}

\section{Study of a dynamic comparator}
\subsection{Architectures}
\subsubsection{Regeneration, reset, and output retention}
\subsubsection{Kickback and its dependence on CDAC state}
\subsection{Input-referred noise}
\subsubsection{Standalone and full-ADC noise simulations}
\subsection{Input offset voltage}
\subsubsection{Offset relative to sampling noise and CDAC error}
\subsection{Clock rate and metastability}
\subsubsection{Evaluation-time and input-overdrive dependence}
\subsubsection{Comparator, logic, and CDAC timing margins}
\subsubsection{Metastability versus incomplete output-latch transfer}
\subsection{Power consumption}
\subsection{Layout area}

\clearpage

\section{Study of a compact CDAC array}
\subsection{Capacitor styles}
\subsubsection{MIM capacitors}
\subsubsection{MOS capacitors}
\subsubsection{POD capacitors}
\subsubsection{MOM capacitors}
\subsubsection{FMOM capacitors}
\subsection{Capacitance density, shielding, and parasitics}
\subsubsection{Connected-layer extraction and capacitance partition}
\subsubsection{Input attenuation and usable voltage range}
\subsection{Capacitance budget and weight architecture}
\subsubsection{Sampling-noise and mismatch constraints}
\subsubsection{Coarse and fine weight implementation}
\subsection{Capacitor mismatch}
\subsubsection{Unit-cell and array-level mismatch models}
\subsubsection{Node-specific mismatch and predicted nonlinearity}
\subsection{Switching scheme}
\subsubsection{Energy, control complexity, and common-mode trajectory}
\subsubsection{Initial CDAC state and comparator common mode}
\subsection{Driver sizing and settling time}
\subsubsection{Capacitive load and bit-dependent settling budget}
\subsection{Power consumption}
\subsection{Layout area}

\clearpage

\section{Study of the complete SAR ADC}
\subsection{CDAC switching scheme}
\subsubsection{Extracted top-plate trajectories and settling}
\subsection{System nonlinearity}
\subsubsection{Transfer gain and full-scale input range}
\subsubsection{Mismatch, parasitics, and INL/DNL simulations}
\subsection{Weight calibration}
\subsubsection{Weight extraction from a ramp and comparator A/B tests}
\subsubsection{Sensitivity to reference and input noise}
\subsection{System sampling rate}
\subsubsection{Conversion-path timing budget}
\subsubsection{Noise and distortion versus sampling rate}
\subsection{System noise and resolution budget}
\subsubsection{Quantization, sampling, comparator, and reference noise}
\subsubsection{Predicted ENOB versus input range and sampling rate}
\clearpage
\subsection{Decision redundancy}
\subsubsection{Error-recovery range of redundant weights}
\subsubsection{Redundant versus minimum-cycle conversion at equal total time}
\subsubsection{Effect of redundancy on noise and decision errors}
\subsection{Control-sequence optimization}
\subsubsection{Sampling-to-first-decision gap}
\subsubsection{Comparator evaluation, reset, and logic duty cycles}
\subsubsection{Continuous conversion and final-decision decoding}
\subsection{Control logic}
\subsubsection{SR-latch and SAR-logic propagation delay}
\subsubsection{Per-channel programmability and clock synchronization}
\subsubsection{Logic contribution to ADC power}
\subsection{System power consumption}
\subsubsection{Block-level and full-ADC simulation comparison}
\subsubsection{Energy per conversion versus sampling rate}

\chapter[FRIDA-1: A first generation 65 nm prototype]{FRIDA-1: A first generation 65\,nm prototype}
\section{Implementation of the building blocks selected in the preceding design chapter}
\subsection{Sampling switch and comparator}
\begin{itemize}
  \item The sampling switch and comparator layouts were manually laid out using a traditional analog-on-top flow. Both were customized from previous designs originating in Cordia.
  \item We used these designs because their architectures and device sizes were already close to the optima found in the preceding optimization steps.
\end{itemize}
\subsection{Capacitor drivers and array}
\begin{itemize}
  \item The drivers were sized from standard cells.
  \item The capacitor array was generated directly from its specification to GDS using a Python script and the KLayout Python API.
\end{itemize}
\subsection{Digital control logic}
\begin{itemize}
  \item The digital control logic was implemented entirely in Verilog. An OpenROAD flow converted it to layout, leaving gaps for the analog blocks.
\end{itemize}

\section{ADC block assembly and channel layout}
\subsection{Assembling the ADC block}
\begin{itemize}
  \item Assemble the individual building blocks into the ADC channel, placing the manually laid-out analog blocks in the gaps left by the digital layout.
  \item Place the capacitor array above the drivers and logic, with a shielding plate between them.
  \item Explain the strategy for isolating the power grids and routing the supplies within the channel.
\end{itemize}
\subsection{Physical verification and completed channel layout}
\begin{itemize}
  \item Check layout connectivity, DRC, LVS, and extraction of the assembled channel.
  \item End with the ADC channel-level layout, showing the building-block placement, capacitor array, shielding, and power grids.
\end{itemize}

\section{Chip-level arrangement and hierarchical assembly of the prototype}
\subsection{Configuration and sequencing}
\begin{itemize}
  \item A synthesized SPI register controls the configuration bits of each ADC.
  \item The decision was made to apply the four ADC sequences externally, so four LVDS inputs were included.
\end{itemize}
\subsection{Analog inputs and comparator-output readout}
\begin{itemize}
  \item The ADC inputs were connected directly to a shared wire bus. The sampling switches of the ADCs choose which ADC is sampling from that bus.
  \item The ADC outputs were not resampled on-chip. Instead, the positive-polarity comparator outputs drive medium-performance LVDS transmitters directly for external measurement.
  \item The SAR logic is still on-chip, but the external DAQ system must acquire each comparator output serially in real time.
\end{itemize}
\subsection{Top-level power grids, pad ring, and assembly}
\begin{itemize}
  \item Explain the chip-level power grids and show a labeled pad ring of the chip.
  \item The top level was assembled hierarchically with OpenROAD, using the ADC block as a hardened macro with a LEF view.
\end{itemize}
\subsection{Top-level layout and fabricated prototype}
\begin{itemize}
  \item End with the top-level layout and an image of the fabricated prototype received in spring 2026.
\end{itemize}

\chapter{FRIDA-1 silicon measurements}

\section{Measurement setup and readout validation}
\subsection{Test board, input drive, references, and clock distribution}
\subsection{DAQ system, FastRX capture, and code decoding}
\subsection{Input-source and differential-amplifier noise}
\subsection{Channel-to-code alignment}
\subsection{Fixed-input tests with and without the amplifier}

\section{Sampling-switch characterization}
\subsection{Track-versus-hold noise and offset separation}
\subsection{Acquisition time and input-dependent settling}

\section{Comparator measurements}
\subsection{Input-referred noise}
\subsubsection{Sampling-duration and clock-phase dependence}
\subsection{Decision errors and metastability}
\subsubsection{Error probability versus input overdrive and evaluation time}
\subsubsection{Reset and output-latch retention checks}
\subsection{Kickback and CDAC-state dependence}
\subsection{Power consumption}

\section{CDAC characterization}
\subsection{Capacitor mismatch}
\subsubsection{Ramp-based and comparator A/B weight extraction}
\subsubsection{Measured weights versus extracted mismatch models}
\subsection{Parasitic capacitance and transfer gain}
\subsection{Initialization state and common-mode trajectory}
\subsection{Power consumption}

\section{Static nonlinearity}
\subsection{Low-noise ramp, full-scale range, and transfer gain}
\subsection{INL and DNL before and after calibration}

\clearpage

\section{Fixed-input noise and channel variation}
\subsection{Code-density distributions across all channels}
\subsection{Sampling-rate dependence of noise and offset}

\section{Calibration methods}
\subsection{Capacitor-weight estimation and repeatability}
\subsection{Noise and linearity after weight correction}

\section{Sampling rate and distortion}
\subsection{ENOB and distortion versus input frequency and rate}
\subsection{Continuous 100 ns operation and acquisition-chain checks}

\section{Timing optimization of control sequence}
\subsection{Comparator and logic clock duty-cycle sweep}
\subsection{Sampling gap, comparator reset, and CDAC-settling margins}
\subsection{Redundant versus minimum-cycle error rates}
\subsection{Synchronized versus independent channel clocks}

\chapter[FRIDA-2: A revised 65 nm design]{FRIDA-2: A revised 65\,nm design}
\section{Sampling and input-path characterization}
\subsection{Sampling-switch behavior with extracted loading}
\subsection{Sampling window and input-noise budget}

\section{Comparator and timing changes}
\subsection{Lower kickback, offset, and metastability risk}
\subsection{Transistor-level layout and extracted verification}
\subsection{Comparator, logic, and CDAC timing allocation}

\section{Capacitor array and driver changes}
\subsection{Connected metal layers and improved parasitic modeling}
\subsection{Capacitor mismatch modeling for linearity studies}
\subsection{Coarse and fine weights and common-mode control}
\subsection{Driver sizing for the extracted capacitance}

\section{Sequence, logic, and system changes}
\subsection{SR-latch and SAR-logic delay reduction}
\subsection{Clock synchronization and channel-to-channel errors}
\subsection{Redundancy versus conversion time and power}

% \chapter{Second-prototype measurements}
% \section{Comparable measurement conditions and readout checks}
% \subsection{Input, reference, and clock conditions matched to the first prototype}
% \subsection{FastRX capture and code-decoding verification}
% \section{Sampling noise, comparator errors, and timing margins}
% \subsection{Track-and-hold noise versus sampling time}
% \subsection{Comparator error rate versus clock phase and overdrive}
% \section{CDAC weights, transfer gain, and static linearity}
% \subsection{Extracted versus measured weights and common mode}
% \subsection{Full-scale range, INL, and DNL}
% \section{Calibration and redundant-decision performance}
% \subsection{Weight calibration and residual nonlinearity}
% \subsection{Redundant versus minimum-cycle conversions}
% \section{ENOB, distortion, and power versus sampling rate}
% \subsection{Fixed-input noise and dynamic-input spectra}
% \subsection{Continuous conversion and energy per sample}
% \section{First- and second-prototype comparison}
% \subsection{Resolution, error rate, and linearity changes}
% \subsection{Power, area, and remaining limitations}

\end{document}
